\documentclass[11pt,a4paper]{article}
\usepackage[margin=2.4cm]{geometry}
\usepackage{booktabs}
\usepackage{array}
\usepackage{xcolor}
\usepackage{longtable}
\usepackage{enumitem}
\usepackage{xurl}
\usepackage{amsmath}
\usepackage[hidelinks,colorlinks=true,linkcolor=black,urlcolor=blue!60!black,citecolor=black]{hyperref}
\usepackage{titlesec}
\usepackage{fancyhdr}

\definecolor{ink}{HTML}{0B0B0B}
\definecolor{muted}{HTML}{52514E}
\definecolor{rule}{HTML}{C9C7C2}

\titleformat{\section}{\large\bfseries}{\thesection}{0.7em}{}
\titleformat{\subsection}{\normalsize\bfseries}{\thesubsection}{0.6em}{}
\setlist[itemize]{leftmargin=1.2em,itemsep=2pt,topsep=3pt}
\setlist[enumerate]{leftmargin=1.4em,itemsep=2pt,topsep=3pt}
\renewcommand{\arraystretch}{1.25}

\pagestyle{fancy}\fancyhf{}
\fancyhead[L]{\footnotesize\color{muted}Weekend Margin Clearinghouse — Pricing and Risk}
\fancyhead[R]{\footnotesize\color{muted}Team 9}
\fancyfoot[C]{\footnotesize\color{muted}\thepage}
\renewcommand{\headrulewidth}{0.4pt}

\newcommand{\note}[1]{{\small\color{muted}#1}}
\newcommand{\bp}{\,\text{bp}}

\begin{document}

\begin{center}
{\LARGE\bfseries Pricing and Risk}\\[4pt]
{\large Where every number comes from, and every case it has to survive}\\[10pt]
{\normalsize FRE-GY 6931 Blockchain \& Cryptocurrency, Fall 2026 — Team 9}\\[2pt]
{\small Sai Kishore (contracts, on-chain) \quad·\quad Swastik Mishra (analysis, product)}\\[2pt]
{\small Prof.\ Richard Radnay \quad·\quad TA Hina Tomar \quad·\quad \today}
\end{center}
\vspace{4pt}

\hrule
\vspace{10pt}

\section*{How to read this document}

There are two classes of number in this project, and keeping them apart is the whole
point of this document.

\textbf{The contract's parameters were chosen.} Eight values are fixed in the deployment,
and five of them set the economics: the 10bp vault discount, the 20\%/10\% initial and
maintenance margin pair, and the 1\%/3\% collateral haircut. They were set by engineering
judgement when the contracts were written, committed before any of the analysis below
existed. Four of the five are immutable on Sepolia; the discount alone can be raised by the owner,
up to a hard cap set at deployment. They are the \emph{subject} of this analysis,
not its output.

\textbf{Everything in this document is derived.} The 6.28bp requirement, the 32.9bp break
point, the 0.456-year duration limit, the 3.98\% typical closure and every figure in
Section~\ref{sec:provenance} come out of the two data series in \S1.1. None was chosen,
and none was adjusted to flatter a parameter.

That ordering — parameters first, measurement second — is the honest one and also the
demanding one: the analysis could have found the chosen values indefensible, and had to be
run in a way that would have shown it. This document does three things with them.

\begin{enumerate}
\item \textbf{Shows where each one came from} — which data series, which formula, which
      judgement call, and which of those judgements is doing the most work.
\item \textbf{Walks every case} — from a median weekend to the worst weekend in eleven
      years of data, in dollars, one step at a time.
\item \textbf{Says where each number breaks} — the exact input value at which it stops
      working, and whether history has already reached that value.
\end{enumerate}

Section~\ref{sec:provenance} is a line-by-line provenance table: every figure quoted
anywhere in our write-up, its source, and the one line of arithmetic that produces it.
Nothing in this document was typed by hand; Appendix~\ref{sec:repro} reproduces all of it
from the cached data files.

Two results are worth stating up front, because they are what the exercise was for.

\textbf{The fee is correctly sized, and we can prove where it stops being correct.} 10bp
against a 6.28bp requirement is $1.59\times$ coverage, and we locate the break point of
every input that feeds it. Two of the four matter; the other two have two to three times
the headroom they need and can be argued about freely without changing a conclusion.
Knowing which is which is the difference between a priced product and a guessed one, and
it is also why one candidate term is deliberately absent: a risk we cannot estimate is
named in \S\ref{sec:unpriced} rather than given an invented number in the fee.

\textbf{And the analysis found the binding constraint, which is not the fee.} Across all
seven historical crises we tested, what threatens the system is the vault running out of
cash — a liquidity limit that sits in the vault's \emph{sizing}, not in its pricing. That
is a more valuable finding than a confirmation would have been: it moves the design
question from ``is 10bp enough?'', which we can now answer yes, to ``how much capital
stands behind it?'', which is where the remaining work belongs. Section~\ref{sec:limits}
develops both.

\section{Inputs}

\subsection{Data}

Four series, cached in \texttt{docs/data/} so every figure is reproducible without a
network call.

\begin{center}
\small
\begin{tabular}{@{}p{3.6cm}p{3.1cm}p{3.0cm}p{4.6cm}@{}}
\toprule
\textbf{Series} & \textbf{Source} & \textbf{Coverage} & \textbf{Used for} \\
\midrule
3-month T-bill yield (DTB3) & FRED, daily, discount basis actual/360 & 2015-01-02 to 2026-10-23, 3{,}085 obs & Weekend yield shocks; the current yield $y$ \\
BTC/USD hourly & Coinbase, OHLC & 2021-01-01 to 2026-10-23, 51{,}070 bars & Closure moves, regime stress, liquidation tests \\
BTC perp funding, hourly & Deribit & matched to the above & Funding-rate realism check \\
Sepolia reconciliation & our own deployment & 28 reads & Confirming the settings below match the deployed contracts \\
\bottomrule
\end{tabular}
\end{center}

\subsection{Settings}

These are not assumptions. Each is read back off the deployed contract in Section~F
of the notebook and matched against the value used here.

\begin{center}
\small
\begin{tabular}{@{}p{5.4cm}p{2.2cm}p{7.4cm}@{}}
\toprule
\textbf{Setting} & \textbf{Value} & \textbf{Where it binds} \\
\midrule
Initial margin & 20\% & Caps opening leverage at $5\times$ \\
Maintenance margin & 10\% & Liquidation trigger \\
Haircut, market open & 1\% & sBILL counted at 99\% of NAV \\
Haircut, market closed & 3\% & sBILL counted at 97\% of NAV \\
sBILL interest & 4.00\% simple & Accrues on wall-clock, weekends included \\
Vault discount & 10bp & Vault buys sBILL at NAV${}\times0.999$. Owner-adjustable up to a 1\% cap, only while the market is open \\
Liquidation reward & 50bp & Paid to the liquidator \\
Funding cap & 10bp/hour & Largest rate the owner can post \\
\bottomrule
\end{tabular}
\end{center}

\subsection{Judgement parameters}

Six inputs are genuine judgement calls. Five of them are currently buried as function
defaults in the notebook rather than declared alongside the settings — a presentation
flaw we name here because these are precisely the numbers a reader should be allowed to
disagree with.

\begin{center}
\small
\begin{tabular}{@{}p{4.5cm}p{1.9cm}p{8.6cm}@{}}
\toprule
\textbf{Parameter} & \textbf{Value} & \textbf{Basis} \\
\midrule
Closure length & 65 hours & Friday 4pm ET to Monday 9am ET. \texttt{MarketCalendar} closes Friday 21:00 to Monday 14:00 UTC, the same 65 hours, so contract and analysis agree. \\
Hurdle premium & 2.00\% & What a vault depositor should want \emph{above} T-bills for taking contract and liquidity risk. Assumed, not observed. \\
Collateral duration & 0.25 yr & Assumes 3-month paper. Section~\ref{sec:breaks} shows this is the tightest assumption in the model. \\
P(long weekend) & 9/52 & Nine of fifty-two closures are extended by a Monday holiday or Good Friday. \\
Extra hours if long & 24 & One additional day of capital lock-up. \\
\bottomrule
\end{tabular}
\end{center}

\section{What actually happens in a closure}

Before the arithmetic, the mechanism. A perpetual swap settles hourly. When the market is
open and a trader's margin falls short, the clearinghouse redeems their sBILL with the
issuer at full NAV. When the market is closed the issuer is unreachable, so the
clearinghouse sells the sBILL to our vault instead, at a discount.

\begin{enumerate}
\item \textbf{Friday 4pm ET.} The T-bill market closes. The clearinghouse switches the
      haircut from 1\% to 3\%, so every sBILL-backed position is now counted slightly
      more conservatively.
\item \textbf{Each hour of the closure.} The price moves. The losing side owes cash and
      holds none, so \texttt{\_collect()} routes to the vault: the vault pays cash, and
      takes sBILL worth slightly more than the cash it paid.
\item \textbf{Monday 9am ET.} The issuer reopens. The vault calls
      \texttt{redeemInventory()}, hands the sBILL back at NAV, and books the difference.
\end{enumerate}

The vault is therefore a weekend \emph{liquidity} provider, not a trader. It earns the
discount, earns the T-bill interest for the hours it holds the paper, and carries the
price risk of that paper across the closure. Those three terms are the whole pricing
problem.

\subsection{The decomposition}

Per \$1 of cash the vault actually deploys, over a closure of $h$ hours:

\[
\underbrace{\frac{(y+\pi)h}{8760}}_{\text{gross carry}}
\;-\;\underbrace{\frac{yh}{8760}}_{\text{sBILL accrual}}
\;+\;\underbrace{D\cdot\Delta y}_{\text{price risk}}
\;+\;\underbrace{p_{\ell}\frac{(y+\pi)h_{\text{extra}}}{8760}}_{\text{long-weekend risk}}
\;=\;\text{required discount}
\]

with $y$ the T-bill yield, $\pi$ the hurdle premium, $D$ duration, $\Delta y$ the weekend
yield shock, and the last term the cost of a Monday holiday extending the closure.

$\Delta y$ is an \textbf{empirical percentile, not a standard deviation}, and the choice is
deliberate. The standard deviation of weekend moves is 4.17bp, but the series is far from
normal: a normal distribution expects 0.04 weekends beyond $4\sigma$ in a sample this size
and we observe nine, and June 2022 is an $8.4\sigma$ move that a normal distribution says
cannot occur. Sigma-scaling would therefore understate the tail by orders of magnitude. We
read the quantile straight off 597 observed weekends instead, which assumes nothing about
the shape of the distribution. The base case uses the 99th percentile, 18bp \mbox{---} 6.5 times
the average absolute move and 4.3 times one standard deviation.

\subsection{What is being priced}

One convention governs every term above, and it is worth stating before the numbers
rather than after them: \textbf{we price the collateral the system is designed for, not
the mock standing in for it.}

Our sBILL accrues a fixed 4\% on a clock, so rates cannot move its value. The dominant
term $D\cdot\Delta y$ therefore prices a risk our own deployment cannot experience, and
that is deliberate: the fee has to be right for the asset the design calls for, not for
the stand-in. Dropping it on the grounds that the mock cannot produce it would leave
$1.48+0.28=1.76$bp against the 10bp charged, and would price a demonstration rather than a
design.

The same convention has a limit, and we draw it at measurability. A risk enters the fee
only when it can be estimated from something. The weekend yield shock is read off 597
observed weekends; the long-weekend term is counted off a calendar, nine closures in
fifty-two. \textbf{Redemption gates meet neither test.} Issuer discretion over redemption
is a real feature of tokenized Treasury products, and our own \texttt{SBill} has a
structural analogue — it pays interest from a reserve funded once at deployment, and
\texttt{redeem()} transfers without checking the balance, so an exhausted reserve reverts
the redemption. But the sector is roughly two years old and has never been stressed, so
any frequency we assigned would be a prior wearing the costume of an estimate. We
therefore \textbf{name this risk and do not price it}, in \S\ref{sec:unpriced}, alongside
the larger unpriced exposure in \S\ref{sec:inversion}. Carrying it in the fee would have
added 0.29bp and changed no conclusion in this document.

Note the first two terms partly cancel: the vault's capital must earn $y+\pi$, but the
sBILL it holds is itself earning $y$, so only the premium $\pi$ has to come out of the
discount. This is why the net carry term is small.

\subsection{Weekend yield shocks}

Measured from each Friday's published DTB3 rate to the next published rate, keeping only
gaps of four days or fewer — 597 weekends.

\begin{center}
\small
\begin{tabular}{@{}lrr@{}}
\toprule
& \textbf{Yield move} & \textbf{Price impact at $D=0.25$} \\
\midrule
Median & 2bp & 0.50bp \\
95th percentile & 10bp & 2.50bp \\
99th percentile & 18.04bp & 4.51bp \\
Largest, 2022-06-10 & 35bp & 8.75bp \\
\bottomrule
\end{tabular}
\end{center}

\note{The four largest weekend moves are 2022-06-10 ($+$35bp), 2023-03-17 ($+$29bp),
2022-07-08 ($+$22bp) and 2022-10-07 ($+$22bp). All four fall in the 2022--23 tightening
cycle, which is the regime this pricing should be judged against.}

\subsection{The three cases}

With $y=4.01\%$ (DTB3, 2026-10-28) and $h=65$:

\begin{center}
\small
\begin{tabular}{@{}lrrrrr@{}}
\toprule
\textbf{Case} & \textbf{Hurdle} & \textbf{Duration} & \textbf{Shock} & \textbf{Terms (bp)} & \textbf{Required} \\
\midrule
Low    & 4.01\% & 1 month  & 10bp & $0.00+0.83+0.19$ & 1.02bp \\
Base   & 6.01\% & 3 months & 18bp & $1.48+4.51+0.28$ & \textbf{6.28bp} \\
Stress & 9.01\% & 3 months & 35bp & $3.71+8.75+0.43$ & 12.89bp \\
\bottomrule
\end{tabular}
\end{center}

The contract charges \textbf{10bp}: $1.59\times$ the base requirement, and comfortably
clear of the low case. The stress case at 12.89bp sits above it, which is the useful
result rather than an awkward one — it means the fee has a locatable limit, and
Section~\ref{sec:breaks} pins that limit to the exact value of each input. A fee that
covered every conceivable case would simply be overpriced.

Observe which term dominates: T-bill \emph{price risk} is 4.51 of the 6.28bp, 72\% of the
total. The single largest component of our fee prices the risk that interest rates move
while the vault is holding the paper.

\section{One trade, five cases}

Take \$100{,}000 of collateral routed through the vault in a single closure. Fee and carry
are fixed; only the rate shock varies.

\begin{center}
\small
\begin{tabular}{@{}lrrrrrr@{}}
\toprule
\textbf{Weekend} & \textbf{Shock} & \textbf{Fee} & \textbf{Carry} & \textbf{Price loss} & \textbf{Kept} & \textbf{vs hurdle \$47.44} \\
\midrule
Typical (median)      & 2.0bp  & \$100.00 & \$29.75 & $-$\$5.00  & \$124.75 & $+$\$77.31 \\
Bad, 1-in-20          & 10.0bp & \$100.00 & \$29.75 & $-$\$25.00 & \$104.75 & $+$\$57.31 \\
Bad, 1-in-100         & 18.04bp & \$100.00 & \$29.75 & $-$\$45.10 & \$84.65  & $+$\$37.21 \\
\textbf{Break point}  & \textbf{32.9bp} & \$100.00 & \$29.75 & $-$\$82.31 & \$47.44 & \textbf{\$0.00} \\
Worst ever (Jun 2022) & 35.0bp & \$100.00 & \$29.75 & $-$\$87.50 & \$42.25  & $-$\$5.19 \\
\bottomrule
\end{tabular}
\end{center}

Each row reads the same way. The vault pays out \$100{,}000 in cash and takes back sBILL
worth \$100{,}100, so the fee is \$100. It holds that paper for 65 hours at 4.01\%, which
is \$29.75. Then rates move, the paper loses $D\times\Delta y$ of its value, and whatever
remains is what the vault keeps. Its depositors required \$47.44 for that weekend.

\subsection{What ``breaks'' means}

The word is doing a lot of work, so it is worth separating two thresholds that are easy to
collapse into one.

\begin{center}
\small
\begin{tabular}{@{}p{4.2cm}p{5.2cm}p{5.6cm}@{}}
\toprule
& \textbf{Missing the hurdle} & \textbf{Losing money} \\
\midrule
What happens & Vault profits, but by less than depositors required & Vault ends the closure with less than it started \\
Threshold, 3-month paper & 32.9bp & 51.9bp \\
Reached in history? & Yes — once, June 2022 at 35bp & No. Largest ever is 35bp \\
Who notices & Depositors earn below hurdle & ERC-4626 share price falls \\
\bottomrule
\end{tabular}
\end{center}

So the correct statement about June 2022 is narrow: on \$100{,}000 the vault came up
\textbf{\$5.19 short} of what it had promised, having still made \$42.25. It is a
knife-edge, not a cliff. Nobody lost principal.

\section{Where each input breaks}
\label{sec:breaks}

Each row solves one input for the value at which 10bp stops clearing the hurdle, holding
the others at base.

\begin{center}
\small
\begin{tabular}{@{}p{4.0cm}p{2.0cm}p{2.4cm}p{2.0cm}p{3.6cm}@{}}
\toprule
\textbf{Input} & \textbf{Base} & \textbf{Breaks above} & \textbf{Headroom} & \textbf{Reached in history?} \\
\midrule
Weekend yield shock & 18bp & \textbf{32.9bp} & $1.83\times$ & \textbf{Yes — 35bp, Jun 2022} \\
Collateral duration & 0.25 yr & \textbf{0.456 yr} & $1.82\times$ & Immediately, on 6-month paper \\
Closure length & 65 h & 228 h & $2.0\times$ vs 113 h & No. Longest closure is 113 h \\
Hurdle premium & 2.0\% & 6.7\% & $3.4\times$ & Not observable \\
\bottomrule
\end{tabular}
\end{center}

Two inputs matter and two do not. The closure length and the hurdle premium have two to
three times the headroom they need, so arguing about them changes nothing. The yield shock
has already been breached once, and the duration assumption has the tightest margin in the
model — and unlike the other three it is an assumption about our own collateral rather
than about the world, which makes it the one we can actually do something about.

\subsection{The duration assumption}

This is the one carrying the most weight, because it is an assumption about the collateral
rather than about the world. Price sensitivity is linear in duration, so doubling the
maturity doubles the loss from any rate move.

\begin{center}
\small
\begin{tabular}{@{}p{4.6cm}rrp{4.4cm}@{}}
\toprule
\textbf{Collateral} & \textbf{Loses above} & \textbf{Worst ever} & \textbf{Verdict} \\
\midrule
3-month bill (our assumption) & 51.9bp & 35bp & $1.48\times$ headroom \\
6-month bill & 26.0bp & 35bp & \textbf{Has happened} \\
1-year note & 13.0bp & 35bp & \textbf{Happens regularly}; 99th pct is already 18bp \\
\bottomrule
\end{tabular}
\end{center}

The design implication is concrete and belongs in the contract rather than the
spreadsheet: \textbf{the vault should accept only short-dated collateral, and that limit
should be enforced, not assumed.} At 3 months the pricing survives the worst weekend on
record. At 6 months it does not survive an ordinary bad one.

\section{Case study: June 2022, the worst weekend}

This is the single most informative closure in the data, because both legs of the risk
fired at once.

On Friday 10 June 2022 the May CPI print landed at 8.6\%, the highest since 1981. Rates
repriced through the close:

\begin{center}
\small
\begin{tabular}{@{}lr@{}}
\toprule
Thursday 9 June & 1.26\% \\
Friday 10 June & 1.34\% \\
Monday 13 June & \textbf{1.69\%} \\
\midrule
Move across the closure & \textbf{$+$35bp} \\
\bottomrule
\end{tabular}
\end{center}

Meanwhile BTC fell from \$29{,}169 at Friday 4pm ET to a low of \$23{,}319 and closed
Monday at \$23{,}610 — a 19.06\% adverse excursion on settled hourly prices, with no
recovery inside the closure.

That combination is what makes it the worst case rather than merely a bad one: \textbf{the
vault buys the most collateral at exactly the moment that collateral is worth the least.}
The two risks are positively correlated, which is the opposite of what a liquidity
provider wants.

On \$1m of notional, a long would have sold roughly \$190{,}567 to the vault:

\begin{center}
\small
\begin{tabular}{@{}lrrrrr@{}}
\toprule
\textbf{Collateral} & \textbf{Fee} & \textbf{Carry} & \textbf{Price loss} & \textbf{Net} & \textbf{Dollars} \\
\midrule
3-month bill & $+$10.00bp & $+$0.99bp & $-$8.75bp & $+$2.24bp & $+$\$42.77 \\
6-month bill & $+$10.00bp & $+$0.99bp & $-$17.50bp & $-$6.51bp & $-$\$123.98 \\
1-year note  & $+$10.00bp & $+$0.99bp & $-$35.00bp & $-$24.01bp & $-$\$457.47 \\
\bottomrule
\end{tabular}
\end{center}

Note the carry: only 0.99bp, against 2.98bp at today's yields. Rates were 1.34\% in 2022,
so the cushion that normally absorbs a shock was mostly absent. \textbf{The pricing is
weakest when rates are low and moving} — which is also when a 35bp weekend becomes
possible. A third correlation working against us.

\section{Regime stress}
\label{sec:regimes}

For each crisis we take the worst closure inside the surrounding window, measured as the
adverse excursion from Friday 4pm ET over 65 hours on settled hourly closes. Excursion is
the right measure rather than the Friday-to-Monday net move, because settlement is hourly:
the losing side pays as the price goes against it, and a later recovery does not return the
collateral already sold.

\begin{center}
\small
\begin{tabular}{@{}p{4.1cm}p{2.0cm}rrrrr@{}}
\toprule
\textbf{Regime} & \textbf{Weekend} & \textbf{BTC open} & \textbf{Excursion} & \textbf{Sold/\$1m} & \textbf{Fee} & \textbf{$\times$typ} \\
\midrule
Yen carry unwind      & 2024-08-02 & \$62{,}553  & 20.38\% & \$203{,}816 & \$204 & $5.1\times$ \\
Jun 2022 CPI shock    & 2022-06-10 & \$29{,}169  & 19.06\% & \$190{,}567 & \$191 & $4.8\times$ \\
SVB / USDC depeg      & 2023-03-10 & \$19{,}972  & 13.54\% & \$135{,}378 & \$135 & $3.4\times$ \\
May 2021 mining ban   & 2021-05-14 & \$49{,}193  & 12.80\% & \$127{,}973 & \$128 & $3.2\times$ \\
LUNA/UST collapse     & 2022-05-06 & \$36{,}081  & 9.10\%  & \$91{,}048  & \$91  & $2.3\times$ \\
Oct 2025 deleveraging & 2025-10-24 & \$110{,}920 & 4.63\%  & \$46{,}333  & \$46  & $1.2\times$ \\
FTX collapse          & 2022-11-11 & \$16{,}573  & 4.45\%  & \$44{,}540  & \$45  & $1.1\times$ \\
\midrule
Typical weekend       & —          & —           & 3.98\%  & \$39{,}795  & \$40  & $1.0\times$ \\
\bottomrule
\end{tabular}
\end{center}

\note{This table measures excursions on hourly closes from Friday 20:00 UTC. The report's
Section~4 table, written by \texttt{docs/stress\_weekends.py}, uses the notebook's convention
(hourly opens, Friday 4pm New York time) and takes the worst closure within three weeks of each
event, which selects 10 October 2025 rather than the 24th. Its figures differ from these by small
amounts (for example 18.33\% rather than 19.06\% for June 2022, and a 95th-percentile capacity of
\$98{,}461 rather than \$93{,}522), and the conclusions below hold under both.}

\subsection{Reading the revenue column honestly}

Crises are good for the vault's revenue line and bad for its risk. Volume scales with the
move, so the yen unwind would have earned $5.1\times$ a typical weekend. But the fee is
10bp of a larger number — \$204 instead of \$40 per \$1m of notional — and nothing about
that \$204 compensates for the two problems the same weekend creates.

\textbf{First, capacity.} Sizing the vault at the 95th percentile of weekend excursions
gives \$93{,}522 of capacity per \$1m of notional. Four of the seven crises exceed it:

\begin{center}
\small
\begin{tabular}{@{}lrr@{}}
\toprule
\textbf{Regime} & \textbf{Needs} & \textbf{\% of capacity} \\
\midrule
Yen carry unwind    & \$203{,}816 & 218\% \\
Jun 2022 CPI shock  & \$190{,}567 & 204\% \\
SVB / USDC depeg    & \$135{,}378 & 145\% \\
May 2021 mining ban & \$127{,}973 & 137\% \\
LUNA/UST collapse   & \$91{,}048  & 97\% \\
Oct 2025            & \$46{,}333  & 50\% \\
FTX collapse        & \$44{,}540  & 48\% \\
\bottomrule
\end{tabular}
\end{center}

A 95th-percentile vault is undersized for exactly the weekends it exists to handle. When
it empties, settlement stops and waits for Monday.

\textbf{Second, liquidation.} Whether the position survives depends entirely on opening
leverage, and the 3\% haircut barely moves the threshold.

\begin{center}
\small
\begin{tabular}{@{}lrrl@{}}
\toprule
\textbf{Leverage} & \textbf{Long liquidated below} & \textbf{Short liquidated above} & \textbf{Crises surviving} \\
\midrule
$2.6\times$ (demo) & $-$31.4\% & $+$25.5\% & all seven \\
$4.0\times$        & $-$16.4\% & $+$13.3\% & four of seven \\
$4.85\times$ (max) & $-$11.5\% & $+$9.3\%  & three of seven \\
\bottomrule
\end{tabular}
\end{center}

At the demo's $2.6\times$ nothing is liquidated in any regime. At $4\times$ the June 2022
and yen-unwind longs go, and the SVB short goes. At the maximum the 2021 mining ban takes
the long too. The demo is safe because it is conservatively levered, not because the
parameters are strong — a distinction worth stating plainly rather than letting the
Sepolia run imply otherwise.

\section{The haircut}

A haircut counts collateral at less than face value, because it is not cash and would have
to be sold. Ours is 1\% when the market is open and 3\% when it is closed. The logic is
right: on a weekday you redeem with the issuer at full NAV, while at a closure your only
buyer is the vault, at a discount — or there is no buyer at all.

\subsection{What it costs}

\begin{center}
\small
\begin{tabular}{@{}lrrr@{}}
\toprule
\textbf{\$30{,}000 posted as} & \textbf{Counts as} & \textbf{Max leverage} & \textbf{Short liquidated at} \\
\midrule
Cash (tUSDC)        & \$30{,}000 & $5.00\times$ & \$98{,}049 ($+$25.94\%) \\
sBILL, 1\% haircut  & \$29{,}700 & $4.95\times$ & \$97{,}959 ($+$25.82\%) \\
sBILL, 3\% haircut  & \$29{,}100 & $4.85\times$ & \$97{,}774 ($+$25.59\%) \\
\bottomrule
\end{tabular}
\end{center}

Moving from the weekday to the weekend haircut brings the liquidation threshold
\textbf{\$185 closer} on a figure near \$98{,}000 — 0.19\% of it — and costs a tenth of a
turn of leverage. The weekend haircut is close to free.

\subsection{What it is sized against}

\begin{center}
\small
\begin{tabular}{@{}lr@{}}
\toprule
Measured T-bill price risk, 99th pct & 4.51bp \\
Worst-ever price risk, 35bp shock & 8.75bp \\
Vault discount paid on conversion & 10.00bp \\
\midrule
Total cost of converting, worst case & 18.75bp \\
\textbf{The haircut} & \textbf{300bp} \\
\midrule
Multiple of measured price risk & $67\times$ \\
Multiple of worst-case conversion cost & $16\times$ \\
\bottomrule
\end{tabular}
\end{center}

So 3\% is one to two orders of magnitude larger than the valuation risk it is usually
described as covering — and measuring that gap is what tells us the haircut is doing a
different job from the one its name suggests.

\textbf{Keep it, and describe it correctly.} It costs \$185 and a tenth of a turn of
leverage, and it is the one place in the contract that becomes more conservative the
moment the market shuts, which is exactly the right instinct. What it buys is not a
calibrated price for 4.5bp of measurable T-bill price risk; it is deliberately generous
cover for the closure risks we cannot put a number on — a longer closure, a delayed
redemption, an issuer problem. At \$185 that cover is close to free, which is the strongest
possible argument for keeping it.

\subsection{What it cannot cover}

The haircut answers ``will this sell below face value?'' — and over a weekend the answer
is yes, by about 10bp. Trivial. The question that matters is ``will it sell at all?''

\begin{center}
\small
\begin{tabular}{@{}lrr@{}}
\toprule
\textbf{BTC move over a 65-hour closure} & \textbf{Size} & \textbf{vs 3\% haircut} \\
\midrule
99th percentile, short side & 13.50\% & $4.5\times$ uncovered \\
99th percentile, long side  & 13.93\% & $4.6\times$ uncovered \\
Worst observed              & 20.38\% & $6.8\times$ uncovered \\
\bottomrule
\end{tabular}
\end{center}

This is not really the haircut's failure. The contract has two distinct buffers, and they
answer different questions:

\begin{itemize}
\item \textbf{The haircut (3\%)} — this asset may not convert at face value. A
      \emph{valuation} adjustment.
\item \textbf{Maintenance margin (10\%)} — the price may move before the position can be
      closed. A \emph{solvency} buffer.
\end{itemize}

A vault with no cash left is a \emph{liquidity} failure, and neither buffer prices it.
Marking the collateral down does not produce a buyer. That risk is carried by vault
sizing, and Section~\ref{sec:regimes} shows the sizing is the weak link.

\section{What this pricing does not cover}
\label{sec:limits}

The analysis establishes three boundaries on its own results. Each one is a finding in
its own right, and each points at a specific change rather than a caveat.

\subsection{The model is deliberately more conservative than the implementation}

\texttt{SBill.nav()} grows purely on the clock:

\begin{center}
\texttt{1e18 + (1e18 * apyBps * (block.timestamp - startTime)) / (10\_000 * 365 days)}
\end{center}

A fixed 4\%, which never reprices when rates move. The 35bp spike of June 2022 would be
invisible to our deployment: the vault would bank its full 10bp and notice nothing.

We priced the risk the collateral would carry in production, not the lighter risk the
mock carries — which is the right way round, and worth being explicit about. T-bill price
risk is 4.51 of the 6.28bp required discount, 72\% of the total, and it is a risk our
deployment \emph{structurally cannot experience}. So the largest component of the fee is
one that no unit test and no Sepolia transaction can exercise. The pricing is built for a
real tokenized T-bill and will hold when one is swapped in; the demo is simply a weaker
test than the model it implements.

\subsection{A loss would be invisible on-chain}

\begin{sloppypar}
If the vault did take a price loss, it surfaces at redemption. The stored
\texttt{inventoryCost} is zeroed, \texttt{totalAssets()} falls, and the ERC-4626
share price drops. Depositors absorb it. Nothing reverts and nothing pauses.
\end{sloppypar}

And the event cannot distinguish a loss from breaking even
(\texttt{WeekendVault.sol:130}):

\begin{center}
\small\texttt{emit InventoryRedeemed(shares, usdcReceived, usdcReceived > cost ? usdcReceived - cost : 0);}
\end{center}

On a loss the profit field emits zero. A depositor reading the log cannot tell ``broke
even'' from ``lost \$457.'' Emitting a signed \texttt{int256}, or emitting \texttt{cost}
alongside \texttt{usdcReceived}, makes losses visible. A one-line fix, and finding a
silent-loss path in one's own code is exactly what a security review is for.

\subsection{Redemption gates: named, not priced}
\label{sec:unpriced}

A tokenized fund can suspend redemptions. If that happened on a Monday the vault would
hold sBILL it could not convert, for an unknown period — and our own \texttt{SBill} has
the same shape of failure, since \texttt{redeem()} transfers from a reserve funded once at
deployment and reverts when that reserve is empty.

We do not price it, and the reason is the measurability test in \S2.2. Every other input
in the fee is anchored to something countable: 597 observed weekends for the yield shock,
fifty-two calendar closures for the long-weekend term. Tokenized Treasury funds are
roughly two years old and have not been through a stress event, so there is no base rate
to estimate from. Any frequency we assigned would be a prior presented as a measurement,
and that is a worse failure than leaving the term out.

Leaving it out is cheap. On the priors we had originally used — a gate once in twenty
years costing 3\% — the term contributes 0.29bp, about 4\% of the fee, and removing it
changes no conclusion in this document: coverage improves from $1.52\times$ to
$1.59\times$ and every break point moves further away. A term that small, resting on a
number that soft, does not belong in a price.

What replaces it is a threshold. The fee stops clearing its hurdle if gates reach roughly
one weekend in eighty, or about once every eighteen months. That is the claim we are
willing to defend — not that gates are rare at some particular rate, but that they are
rarer than once every eighteen months, which the existence of a \$15bn sector makes
difficult to dispute. The exposure is managed by diversifying issuers and by sizing the
vault, not by charging for it.

\subsection{The inversion}
\label{sec:inversion}

The largest unpriced exposure is the one the vault holds all the time.

\begin{center}
\small
\begin{tabular}{@{}lr@{}}
\toprule
Priced: T-bill price risk & 4.51bp \\
Priced: net carry & 1.48bp \\
Priced: long-weekend risk & 0.28bp \\
\midrule
\textbf{Total risk budget} & \textbf{6.28bp} \\
\midrule
\textbf{Not priced:} redemption gate (\S\ref{sec:unpriced}) & no base rate \\
\textbf{Not priced:} USDC depeg, March 2023 & \textbf{1{,}300bp} \\
\quad as a multiple of the entire budget & $207\times$ \\
\quad as years of fee revenue, \$1m book & 89.5 years \\
\bottomrule
\end{tabular}
\end{center}

The model prices four and a half basis points of risk on paper the vault holds for 65
hours, and zero on the stablecoin it holds permanently. In March 2023 USDC traded at
\$0.87. That is 1{,}300bp — more than two hundred times the whole risk budget.

It is also wrong-way. The vault holds the most cash during a crisis, because that is when
settlement draws on it hardest — and a crisis is when stablecoins depeg. The SVB weekend
appears twice in this document for that reason: as the third-largest draw in
Section~\ref{sec:regimes}, and as the largest unpriced loss here. The same weekend
maximises both. Circle's own stated trigger was that its wires had not settled before the
weekend, which is the identical mechanism our system is built to address.

\subsection{What we would change}

\begin{enumerate}
\item \textbf{Cap collateral duration in the contract.} The tightest assumption in the
      model, and currently unenforced.
\item \textbf{Cap or reserve against the stablecoin leg.} We do not claim this exposure
      can be priced into the discount — a 1{,}300bp tail against a 6.28bp budget is not a
      charge any fee absorbs. It is a position limit: cap the vault's idle balance, hold
      part of it in a second stablecoin, or accept the exposure explicitly rather than
      silently.
\item \textbf{Size the vault above the 95th percentile.} Four of seven historical crises
      exhaust it.
\item \textbf{Emit signed profit and loss} from \texttt{redeemInventory()}.
\item \textbf{Promote the five hidden judgement parameters} into the settings cell, so a
      reader can disagree with them.
\end{enumerate}

\section{Figure provenance}
\label{sec:provenance}

Every number quoted in our write-up, with its source and the arithmetic that produces it.

\begin{center}
\footnotesize
\begin{longtable}{@{}p{4.4cm}p{1.8cm}p{3.1cm}p{5.4cm}@{}}
\toprule
\textbf{Figure} & \textbf{Value} & \textbf{Source} & \textbf{How computed} \\
\midrule
\endfirsthead
\toprule
\textbf{Figure} & \textbf{Value} & \textbf{Source} & \textbf{How computed} \\
\midrule
\endhead
Current 3m yield & 4.01\% & FRED DTB3, 2026-10-28 & Last non-null observation \\
Weekend shock, median & 2bp & FRED DTB3 & Median $|r_{t+1}-r_t|$ over 597 Fri-to-next pairs \\
Weekend shock, 99th & 18bp & FRED DTB3 & 99th pct of the same \\
Weekend shock, max & 35bp & FRED DTB3, 2022-06-10 & Max of the same \\
Hurdle per closure & 4.74bp & Derived & $(y+\pi)\tfrac{65}{8760}+p_\ell(y+\pi)\tfrac{24}{8760}$ \\
Required discount, base & 6.28bp & Derived & $1.48+4.51+0.28$ \\
Charged discount & 10bp & \texttt{WeekendVault} & \texttt{discountBps}, read on Sepolia \\
Coverage ratio & $1.59\times$ & Derived & $10/6.28$ \\
Carry on \$100k, 65h & \$29.75 & Derived & $100{,}000\times0.0401\times65/8760$ \\
Break point, yield shock & 32.9bp & Derived & Solve $10\bp+\text{carry}-0.25\Delta y=4.74\bp$ \\
Break point, duration & 0.456 yr & Derived & Solve the same for $D$ at $\Delta y=18$bp \\
Principal loss, 3-month & 51.9bp & Derived & Solve net kept ${}=0$ at $D=0.25$ \\
Principal loss, 6-month & 26.0bp & Derived & Same at $D=0.5$ \\
Jun 2022 BTC excursion & 19.06\% & Coinbase hourly & $1-\min(\text{close})/\text{close}_{\text{Fri 20:00}}$, 65h \\
Jun 2022 net, 3-month & $+$2.24bp & Derived & $10+0.99-8.75$, carry at $y=1.34\%$ \\
Typical weekend & 3.98\% & Coinbase hourly & Mean worst-side excursion, 304 weekends \\
Worst weekend & 20.38\% & Coinbase, 2024-08-02 & Max of the same \\
99th pct, long side & 13.93\% & Coinbase hourly & 99th pct of down-excursion \\
Vault capacity, 95th pct & \$93{,}522 & Coinbase hourly & 95th pct excursion $\times$ \$1m \\
Haircut, counts as & \$29{,}100 & Derived & $30{,}000\times(1-0.03)$ \\
Liquidation, weekend & \$97{,}774 & Derived & $h(D+P_0)/(m+h)$, $h=0.97$, $m=0.10$ \\
Haircut cost & \$185 & Derived & $97{,}959-97{,}774$ \\
Haircut vs price risk & $67\times$ & Derived & $300/4.51$ \\
USDC depeg & 1{,}300bp & Mar 2023, \$0.87 low & $(1-0.87)\times10{,}000$ \\
Depeg vs budget & $207\times$ & Derived & $1300/6.28$ \\
\bottomrule
\end{longtable}
\end{center}

\appendix
\section{Reproducing these figures}
\label{sec:repro}

All four data series are cached in \texttt{docs/data/}, so no network access is needed.

\begin{enumerate}
\item \texttt{docs/analysis.ipynb} Sections A--E produce the pricing decomposition,
      sensitivity grids and closure statistics. Section~F reads the deployed Sepolia
      contracts and reconciles the settings in \S1.2 against them; it is the only section
      that touches the network.
\item The settings in \S1.2 are mirrored from the deploy script
      \texttt{00\_deploy\_weekend\_margin.ts} in
      \texttt{packages/hardhat/deploy/}. They are not re-typed here: Section~F
      asserts equality and fails loudly otherwise.
\item Figures in this document were recomputed independently of the notebook, from the
      same cached CSVs, as a cross-check. Two discrepancies surfaced and were resolved in
      favour of the stricter definition: closure moves are measured as \emph{adverse
      excursion} rather than Friday-to-Monday net move, because hourly settlement draws on
      the vault at the worst point of the closure and a later recovery does not return
      collateral already sold; and regime weekends are selected as the worst closure within
      a dated crisis window rather than hand-picked, so the choice is reproducible.
\end{enumerate}

\section{Scope of the analysis}

\begin{itemize}
\item \textbf{sBILL is a mock.} It accrues a fixed 4\% from a pre-funded reserve. It
      demonstrates the integration; it does not carry real T-bill exposure, and as
      \S\ref{sec:limits} sets out it cannot express the price risk that dominates the fee.
\item \textbf{Day-count mismatch.} DTB3 is quoted on a discount basis, actual/360. sBILL
      accrues simple interest on 365 days. The difference is immaterial at 65-hour
      horizons but is a real inconsistency.
\item \textbf{BTC history starts in 2021.} 304 weekends. The yield series reaches back to
      2015 and 597 weekends, so the two samples are not co-extensive; the regime table is
      bounded by the shorter one.
\item \textbf{Hourly closes, not trades.} Excursions are measured on settled hourly closes,
      matching how the contract settles. Intra-hour highs and lows are worse — the 99th
      percentile long excursion rises from 13.93\% to 17.42\% — so a liquidation engine
      reading a live price would fire more often than this analysis implies.
\item \textbf{No holiday calendar.} \texttt{MarketCalendar} tests time-of-week only. On a
      trading holiday the contract believes the market is open, routes to a closed issuer,
      locks the vault out, and applies the gentle 1\% haircut. With real collateral
      \texttt{redeem()} would revert and take \texttt{settle()} with it. The 113-hour
      Thanksgiving closure is priced in \S\ref{sec:breaks} but not implemented.
\item \textbf{The hurdle premium is assumed.} 2\% above T-bills is a judgement with no
      observable market price. It has $3.4\times$ headroom, so it is not load-bearing.
\item \textbf{Redemption gates are not priced.} The sector has no base rate, so the fee
      excludes them and \S\ref{sec:unpriced} states the threshold at which they would
      matter instead. A reader who believes gates are more frequent than one weekend in
      eighty should treat the fee as uncovered for that risk.
\end{itemize}

\end{document}
